% Ad Atticum 7.5 --- headnote
% ad-atticum-07-05, 18 December 50 BC, Formiae

\textit{Cicero to Atticus, written from his Formian villa on
the fifteenth day before the Kalends of January, 18 December
50 BC (Perseus dateline:
\textit{Scr. in Formiano xv K. Ian. a. 704 (50)}). The first
of an unbroken sequence of Formian letters in the closing
weeks of the year: Cicero has landed in Italy, set himself up
at the suburban villa just off the Appian Way, and is feeding
Atticus running bulletins as the crisis with Caesar comes to
its head. He is still waiting on a supplicatio for his
Cilician campaign, and a hoped-for triumph, and is keeping
out of the City until matters at Rome clarify; meanwhile he
listens to whatever rumor walks in the door.}

\textit{The substantive content is in sections 4 and 5. The
upper sections are domestic --- worry about Atticus's and
Pilia's quartan fever, and about his freedman Tiro's health
(``it is for his humanity and modesty that I would rather
have him safe than for my own use of him''); a travel
itinerary from Formiae over the Pomptine marshes to Pompey's
Alban villa and so to Rome by his birthday, the 3rd of
January. Then the lid comes off: ``About the commonwealth I
grow more afraid by the day. For the loyalists, as people
suppose them, are not of one mind. What Roman knights, what
senators I have seen, who poured abuse on Pompey for
everything else, but above all for this journey of his!
Peace is what we need. From victory there will come many
evils and surely a tyrant.'' Section 5 then closes the
position. ``I am of the school which holds it more useful
to concede what he demands than to join battle. We resist
too late one whom for ten years we have nourished against
ourselves.'' The Greek leave-taking, \emph{apotripsai},
``shake off,'' the quartan ague --- restores the affectionate
tone the politics had broken.}
